%%%%%%%%%%%%%%%%%%%%%%%%%%%%%%%%%%%%%%%%%%%%%%%%%%%%%%%%%
% Relatorio do 1o experimento - dados simulados
% ATENCAO: substitua as medicoes e os links por registros reais
% antes de entregar o relatorio.
%%%%%%%%%%%%%%%%%%%%%%%%%%%%%%%%%%%%%%%%%%%%%%%%%%%%%%%%%

\documentclass[12pt]{article}

% Usa o modelo SBC quando ele estiver disponível no projeto. As definições
% alternativas abaixo permitem compilar também sem o arquivo sbc-template.sty.
\IfFileExists{sbc-template.sty}{
    \usepackage{sbc-template}
}{
    \newcommand{\email}[1]{\texttt{##1}}
    \newcommand{\address}[1]{\date{##1}}
    \newenvironment{resumo}
        {\renewcommand{\abstractname}{Resumo}\begin{abstract}}
        {\end{abstract}}
}
\usepackage[utf8]{inputenc}
\usepackage[T1]{fontenc}
\IfFileExists{brazil.ldf}{\usepackage[brazil]{babel}}{\usepackage[english]{babel}}
\usepackage{graphicx}
\usepackage{float}
\usepackage{amsmath}
\usepackage{circuitikz}
\usepackage{hyperref}

\hypersetup{
    colorlinks=true,
    linkcolor=black,
    urlcolor=blue,
    citecolor=black
}

\sloppy

\title{Experimento 1\\
Portas Lógicas AND, OR e NOT}

\author{Grupo [número do grupo]\\
        Allan Victor Almeida Faria, 261010546\\
        Caio César Gonçalves Ribeiro, 211038217\\
        \mbox{[Nome do integrante, matrícula]}}

\address{Departamento de Ciência da Computação -- Universidade de Brasília (UnB)\\
  CIC0231 -- Laboratório de Circuitos Lógicos\\
  \today
  \email{[inserir e-mails do grupo]}
}

\begin{document}

\maketitle

\renewcommand{\abstractname}{Abstract}
\begin{abstract}
This report describes a simulated version of the first experiment with AND,
OR, and NOT logic gates. The operation of the digital kit, the output voltages
of the 74HC08 and 74HC32 integrated circuits, and the implementation of an AND
gate using only OR and NOT gates were analyzed. The simulated results reproduced
the expected truth tables. The values and media links must be replaced by actual
experimental records before submission.
\end{abstract}

\begin{resumo}
Este relatório apresenta uma versão simulada do primeiro experimento com portas
lógicas AND, OR e NOT. Foram analisados o funcionamento do kit digital, as
tensões de saída dos circuitos integrados 74HC08 e 74HC32 e a implementação de
uma porta AND usando somente portas OR e NOT. Os resultados simulados reproduziram
as tabelas verdade esperadas. Os valores e links devem ser substituídos pelos
registros reais antes da entrega.
\end{resumo}

\section{Introdução}
\label{sec:introducao}

Os circuitos digitais representam informações por meio de níveis discretos de
tensão, interpretados como valores lógicos 0 e 1 \cite{wakerly90}. As operações fundamentais da
álgebra booleana incluem AND, OR e NOT. Essas operações podem ser implementadas
por circuitos integrados que reúnem várias portas lógicas em um único componente.

Neste experimento foram utilizados os circuitos integrados 74HC08, formado por
quatro portas AND de duas entradas; 74HC32, formado por quatro portas OR de duas
entradas; e 74HC04, composto por portas inversoras. Também foi analisada a
diferença entre os valores lógicos abstratos e as tensões elétricas efetivamente
observadas nas saídas.

\subsection{Objetivos}
\label{sec:objetivos}

Os objetivos do experimento foram conhecer o funcionamento básico do kit lógico,
medir a tensão de alimentação, verificar experimentalmente as tabelas verdade das
portas AND e OR e implementar uma porta AND utilizando apenas portas OR e NOT,
com base nas leis de De Morgan.

\subsection{Materiais}
\label{sec:materiais}

Foram considerados os seguintes materiais e equipamentos:

\begin{itemize}
    \item kit lógico digital com chaves e LEDs;
    \item \textit{protoboard};
    \item multímetro digital;
    \item fios conectores;
    \item circuito integrado 74HC08;
    \item circuito integrado 74HC32;
    \item circuito integrado 74HC04.
\end{itemize}

\section{Procedimentos e Resultados}
\label{sec:procedimentos}

\textbf{Aviso:} as medições e os links apresentados nesta versão são simulados
e servem apenas para estruturar e testar o documento. Eles não constituem
evidência da execução do experimento.

\subsection{Teste do LED e medição da alimentação}
\label{sec:alimentacao}

Inicialmente, um LED do painel digital foi conectado a uma das chaves do kit.
No comportamento simulado, o LED acendeu quando a chave foi colocada em nível
lógico 1 e apagou quando ela foi colocada em nível lógico 0.

Considerou-se a medição da tensão entre os terminais VCC e GND com o multímetro
ajustado para tensão contínua. O valor simulado foi:

\begin{equation}
    V_{CC}=4{,}97\ \text{V}.
\end{equation}

Esse valor está próximo da alimentação nominal de 5 V indicada para o kit.

\subsection{Teste das portas AND e OR}
\label{sec:portas}

As chaves A e B foram consideradas como entradas dos circuitos integrados. As
saídas das quatro portas internas foram associadas aos LEDs L0, L1, L2 e L3.
Para cada combinação das entradas, foram geradas tensões simuladas compatíveis
com o comportamento lógico esperado.

O registro fotográfico fictício da montagem pode ser acessado neste
\href{https://example.com/experimento1/foto-montagem-74hc}{link da foto simulada}.

\subsubsection{Portas AND do CI 74HC08}

Cada porta do CI 74HC08 realiza a operação:

\begin{equation}
    S=A\cdot B.
\end{equation}

A Tabela~\ref{tab:and} apresenta as tensões simuladas para as quatro portas AND.

\begin{table}[H]
    \centering
    \caption{Tensões simuladas nas saídas das portas AND do CI 74HC08.}
    \label{tab:and}
    \begin{tabular}{|c|c|c|c|c|c|}
        \hline
        B & A & L3 = S4 (V) & L2 = S3 (V) & L1 = S2 (V) & L0 = S1 (V) \\
        \hline
        0 & 0 & 0,03 & 0,04 & 0,02 & 0,05 \\
        \hline
        0 & 1 & 0,05 & 0,03 & 0,04 & 0,06 \\
        \hline
        1 & 0 & 0,04 & 0,06 & 0,03 & 0,05 \\
        \hline
        1 & 1 & 4,91 & 4,89 & 4,93 & 4,90 \\
        \hline
    \end{tabular}
\end{table}

Os dados simulados reproduzem a tabela verdade da operação AND. Tensões entre
0,02 V e 0,06 V foram geradas quando ao menos uma entrada estava em nível baixo.
Quando A e B estavam simultaneamente em nível alto, as saídas ficaram entre
4,89 V e 4,93 V.

\subsubsection{Portas OR do CI 74HC32}

Cada porta do CI 74HC32 realiza a operação:

\begin{equation}
    S=A+B.
\end{equation}

A Tabela~\ref{tab:or} apresenta as tensões simuladas para as quatro portas OR.

\begin{table}[H]
    \centering
    \caption{Tensões simuladas nas saídas das portas OR do CI 74HC32.}
    \label{tab:or}
    \begin{tabular}{|c|c|c|c|c|c|}
        \hline
        B & A & L3 = S4 (V) & L2 = S3 (V) & L1 = S2 (V) & L0 = S1 (V) \\
        \hline
        0 & 0 & 0,04 & 0,03 & 0,06 & 0,05 \\
        \hline
        0 & 1 & 4,90 & 4,92 & 4,88 & 4,91 \\
        \hline
        1 & 0 & 4,89 & 4,93 & 4,91 & 4,90 \\
        \hline
        1 & 1 & 4,92 & 4,90 & 4,94 & 4,91 \\
        \hline
    \end{tabular}
\end{table}

Os dados simulados também reproduzem a tabela verdade da porta OR. Para
$A=0$ e $B=0$, as tensões ficaram próximas de 0 V. Nas demais combinações,
as tensões permaneceram próximas da alimentação do circuito, representando o
nível lógico alto.

O procedimento fictício de teste das portas pode ser acessado neste
\href{https://example.com/experimento1/video-teste-and-or}{link do vídeo simulado}.

\subsubsection{Análise das tensões}

Embora os sinais sejam interpretados como valores lógicos 0 ou 1, suas tensões
reais não precisam ser exatamente iguais a 0 V ou 5 V. Usando como referência
pedagógica as faixas apresentadas na Figura 2 do roteiro, as tensões simuladas
próximas de 0 V representam nível lógico baixo, enquanto as tensões próximas de
5 V representam nível lógico alto. Nenhum valor simulado permaneceu na região
de indefinição.

As pequenas diferenças entre S1, S2, S3 e S4 representam possíveis variações
internas das portas, precisão do multímetro, resistência dos fios e contatos,
variação da alimentação e carga introduzida pelos LEDs. Apesar dessas diferenças,
todas as saídas conservaram os níveis lógicos esperados.

Cabe observar que a Figura 2 do roteiro apresenta níveis característicos da
família TTL, enquanto os circuitos utilizados pertencem à família CMOS 74HC.
Por isso, uma análise elétrica rigorosa deve considerar os limites especificados
nos \textit{datasheets} dos componentes 74HC08 e 74HC32.

\subsection{Implementação de uma porta AND com OR e NOT}
\label{sec:and-or-not}

A porta AND pode ser implementada utilizando somente portas OR e NOT por meio
da lei de De Morgan \cite{mortari:16}:

\begin{equation}
    A\cdot B=\overline{\overline{A}+\overline{B}}.
    \label{eq:demorgan}
\end{equation}

As entradas A e B são inicialmente invertidas por duas portas do CI 74HC04.
Os sinais resultantes são aplicados a uma porta OR do CI 74HC32. A saída da
porta OR passa por uma terceira porta NOT, produzindo a operação AND.

\begin{figure}[H]
    \centering
    \begin{circuitikz}
        \node[not port] (notA) at (0,1.5) {};
        \node[not port] (notB) at (0,0) {};
        \node[or port]  (or1)  at (3,0.75) {};
        \node[not port] (notS) at (5.5,0.75) {};

        \draw (notA.out) -- (or1.in 1);
        \draw (notB.out) -- (or1.in 2);
        \draw (or1.out) -- (notS.in);

        \draw (notA.in) -- ++(-1,0) node[left] {$A$};
        \draw (notB.in) -- ++(-1,0) node[left] {$B$};
        \draw (notS.out) -- ++(1,0) node[right] {$S=A\cdot B$};
    \end{circuitikz}
    \caption{Implementação de uma porta AND utilizando portas OR e NOT.}
    \label{fig:and-or-not}
\end{figure}

\begin{table}[H]
    \centering
    \caption{Resultados simulados da porta AND implementada com OR e NOT.}
    \label{tab:and-or-not}
    \begin{tabular}{|c|c|c|c|}
        \hline
        B & A & Saída esperada & Tensão simulada em L0 (V) \\
        \hline
        0 & 0 & 0 & 0,07 \\
        \hline
        0 & 1 & 0 & 0,05 \\
        \hline
        1 & 0 & 0 & 0,06 \\
        \hline
        1 & 1 & 1 & 4,87 \\
        \hline
    \end{tabular}
\end{table}

Os resultados simulados mostram nível lógico alto somente quando A e B estão
simultaneamente em nível 1. Assim, a montagem reproduz a tabela verdade da porta
AND e ilustra a aplicação da lei de De Morgan.

O registro fictício dessa montagem pode ser acessado neste
\href{https://example.com/experimento1/video-and-com-or-not}{link do vídeo simulado}.

\section{Conclusões}
\label{sec:conclusao}

Com base nos dados simulados, foi possível representar o funcionamento das
portas AND e OR e distinguir os níveis lógicos das tensões físicas. As quatro
portas de cada circuito integrado apresentaram o mesmo comportamento lógico,
apesar das pequenas variações de tensão geradas entre as saídas.

A implementação da porta AND com portas OR e NOT reproduziu a tabela verdade
esperada e demonstrou a aplicação prática da lei de De Morgan. Para que estas
conclusões representem a execução efetiva do experimento, as medições simuladas
e os links fictícios devem ser substituídos pelos registros obtidos no laboratório.

\newpage
\section*{Autoavaliação}

\begin{enumerate}
    \item b
    \item a,c,d
    \item \begin{table} [H]
    \centering
    \caption{Tabela Verdade do Exercício 3}
    \label{tab:tabela-autoav}
    \begin{tabular}{|c|c|c|c|}
        \hline
        A & B & C & \(\overline{AB} + C\)\\
        \hline
        0 & 0 & 0 & 1 \\
        \hline
        0 & 0 & 1 & 1 \\
        \hline
        0 & 1 & 0 & 1 \\
        \hline
        0 & 1 & 1 & 1 \\
        \hline
        1 & 0 & 0 & 1 \\
        \hline
        1 & 0 & 1 & 1 \\
        \hline
        1 & 1 & 0 & 0 \\
        \hline
        1 & 1 & 1 & 1 \\
        \hline
    \end{tabular}        
    \end{table}
    \item b,d,a,c (De cima pra baixo)
\end{enumerate}

\bibliographystyle{plain}
\bibliography{relatorio}

\end{document}
